\subsection*{Data.} The empirical foundation of this study was a dataset composed of 424 United States-listed equities and exchange-traded funds spanning 10 years (2014-2024). A pipeline was developed to support the automated construction of excess growth rate models for any ticker within this dataset, allowing for batch simulation and stress testing across market sectors. To validate the performance of this pipeline, we focused our analysis on the single-asset SPY, which tracked the Standard \& Poor's (S\&P) 500 index. The data included daily open, high, low, and close prices along with volume-weighted average price metrics (VWAP). For training purposes, the first 2,766 observations constituted the in-sample dataset, spanning from January 3, 2014, to December 31, 2024. A subsequent 249 trading days, from January 2 to December 31, 2025, were reserved for out-of-sample testing to assess model generalizability.

\subsection*{Excess growth rate calculation.}
We computed the excess growth rate $G_{i,j}$ for ticker $i$ between time period $j-1\rightarrow{j}$ for a given equity price series $P_{i,j}$ (units: dollars per share (USD/share)). To measure how much each stock's growth exceeded a risk-free baseline, we defined the excess growth rate as:
\begin{equation}
    G_{i,j} \equiv \left( \frac{1}{\Delta t} \right) \cdot \ln \left( \frac{P_{i,j}}{P_{i,j-1}} \right) - r_f
    \label{eq:growth_rate}
\end{equation}
where $\Delta t$ represented the time step set to $1/252$ for daily data (units: years), and $r_f$ denoted the constant continuously compounded risk-free rate (units: year$^{-1}$) derived from Separate Trading of Registered Interest and Principal of Securities (STRIPS) bond yields observed on September 10, 2025. Unlike a simple return, the excess growth rate has units of inverse time (year$^{-1}$) and directly quantifies the premium earned above a continuously compounded risk-free benchmark; log growth rates are also time-additive, making them consistent with a continuous compounding framework. While $r_f$ varied dynamically over the sample period, we used a constant proxy to isolate the volatility patterns in the data without adding noise from interest rate movements.

\subsection*{Discrete hidden Markov model.}
We defined the hidden Markov model using the tuple $\mathcal{M} = (\mathcal{S}, \mathcal{O}, \mathbf{T}, \mathbf{E}, \bar{\pi})$ \cite{rabiner_introduction_1986}. The hidden state $S_t \in \mathcal{S} = \{1,\dots,N\}$ represented market mood regimes and was defined by quantile bins of a fitted cumulative distribution function (CDF) for the excess growth rate series. The observation space $\mathcal{O} \subset \mathbb{R}$ consisted of continuous excess growth rate values, the transition matrix $\mathbf{T}$ governed regime-to-regime dynamics, the emission model $\mathbf{E}$ was parameterized by state-conditional Normal distributions, and $\bar{\pi}$ denoted the stationary distribution. The state-conditional observation model specified that the excess growth rate, given hidden state $k$, followed a Normal distribution:
\begin{equation}
    G_t \mid S_t = k \sim \mathcal{N}(\mu_k, \sigma_k).
\end{equation}
To map observations to hidden states during estimation, we defined quantile boundaries $Q_k = F_L^{-1}(k/N; \mu_L, b_L)$ for $k \in \{1,\dots,N-1\}$, with $Q_0 = -\infty$ and $Q_N = +\infty$ so that the partition covers the full support of the distribution. An observation was assigned to state $k$ when $Q_{k-1} < G_t \le Q_k$.

The choice of the Laplace distribution over alternatives such as the Student's t-distribution was motivated by its alignment with the peaks and heavy tails of high-frequency financial excess growth rates \cite{cont_empirical_2001, kotz_laplace_2001}. The Student's t-distribution is often used to model leptokurtosis, but the Laplace distribution's sharper peak at the mean better matched the concentration of small price movements observed in the SPY dataset \cite{mandelbrot_variation_1963, harwood_against_2019}. The Laplace distribution also provided a closed-form expression for quantile estimation, whereas MLE for the Student's t-distribution required iterative numerical optimization \cite{bilmes_gentle_nodate}; this computational simplicity supported the 424-ticker pipeline. The quantile boundaries were derived from a parametric Laplace fit and therefore did not enforce equal counts per bin in the historical data; instead, they provided a partition of the distribution that preserved tail shape. By using the Laplace distribution to define the bins, the model established a baseline characterization of the return distribution's shape before the transition dynamics and jump mechanism were layered on top \cite{kou_jump-diffusion_2002}.

We estimated the transition matrix $\mathbf{T}$ through direct frequentist counting rather than iterative optimization. This approach differed from the traditional expectation-maximization (EM) algorithm used for parameter estimation in Gaussian mixture and hidden Markov models \cite{bilmes_gentle_nodate}. However, the discrete state assignments allowed for straightforward counting of observed transitions between regimes, making the approach both conceptually simple and free of initialization sensitivity. A consequence of direct counting was that empirical transition matrices often had low probabilities for exiting central states into tail states, which caused models to lose the effect of volatility clustering too quickly \cite{bulla_stylized_2006}. We corrected this by adding jump parameters $\epsilon$, representing the probability of a jump event, and $\lambda$, representing the mean duration, along with a tail-state selection rule. Tail-states were defined as $\mathcal{S}_{bottom} = \{1, \dots, N_{tail}\}$ and $\mathcal{S}_{top} = \{N-N_{tail}+1, \dots, N\}$, where $N_{tail}$ was user-configurable. If a jump was triggered, the jump-duration $K$ was sampled from a Poisson distribution such that $K \sim \text{Poisson}(\lambda)$, which was standard for modeling independent events in fixed intervals \cite{glasserman_monte_2003}. During these $K$ steps, the standard Markovian transitions were overridden to target tail-states, with a user-configurable bias $p_{neg}$ (default 0.52) that favored negative-tail states to reproduce gain/loss asymmetry.

% ── Figure 2: Model Architecture ─────────────────────────────────────────────
\begin{figure}[tp]
\centering
\resizebox{\textwidth}{!}{%
\begin{tikzpicture}[
  % ----- node styles -----
  state/.style    = {circle, draw=black, very thick, minimum size=1.1cm,
                     inner sep=0pt, font=\small},
  tailbot/.style  = {state, fill=red!25,  draw=red!70!black},
  tailtop/.style  = {state, fill=blue!25, draw=blue!70!black},
  midstate/.style = {state, fill=white,   draw=black!70},
  % ----- arrow styles -----
  markov/.style   = {<->, >=Stealth, thin, gray!80},
  trigger/.style  = {->, >=Stealth, thick, dashed, black!80},
  jumparc/.style  = {->, >=Stealth, very thick},
]

  % ===== State nodes along the horizontal axis =====
  \node[tailbot]              (s1)   at ( 0.0, 0) {$1$};
  \node[tailbot]              (s2)   at ( 2.0, 0) {$2$};
  \node[font=\Large]          (ld)   at ( 4.0, 0) {$\cdots$};
  \node[midstate]             (sk)   at ( 6.0, 0) {$k$};
  \node[font=\Large]          (rd)   at ( 8.0, 0) {$\cdots$};
  \node[tailtop]              (sNm1) at (10.0, 0) {$N\!-\!1$};
  \node[tailtop]              (sN)   at (12.0, 0) {$N$};

  % ===== Markov transition arrows (bidirectional, thin) =====
  \draw[markov] (s1)  -- (s2);
  \draw[markov] (s2)  -- (ld);
  \draw[markov] (ld)  -- (sk);
  \draw[markov] (sk)  -- (rd);
  \draw[markov] (rd)  -- (sNm1);
  \draw[markov] (sNm1)-- (sN);

  % ===== Laplace emission bell curves above each visible state =====
  \draw[red!75!black, thick]
    plot[domain=-0.42:0.42, samples=40, smooth, variable=\t]
    ({ 0.0+\t}, {1.45 + 0.72*exp(-\t*\t/0.024)});
  \draw[red!25, thin] (-0.42,1.45) -- (0.42,1.45);

  \draw[red!75!black, thick]
    plot[domain=-0.42:0.42, samples=40, smooth, variable=\t]
    ({ 2.0+\t}, {1.45 + 0.72*exp(-\t*\t/0.024)});
  \draw[red!25, thin] (1.58,1.45) -- (2.42,1.45);

  \draw[black!55, thick]
    plot[domain=-0.42:0.42, samples=40, smooth, variable=\t]
    ({ 6.0+\t}, {1.45 + 0.72*exp(-\t*\t/0.024)});
  \draw[gray!40, thin] (5.58,1.45) -- (6.42,1.45);

  \draw[blue!75!black, thick]
    plot[domain=-0.42:0.42, samples=40, smooth, variable=\t]
    ({10.0+\t}, {1.45 + 0.72*exp(-\t*\t/0.024)});
  \draw[blue!25, thin] (9.58,1.45) -- (10.42,1.45);

  \draw[blue!75!black, thick]
    plot[domain=-0.42:0.42, samples=40, smooth, variable=\t]
    ({12.0+\t}, {1.45 + 0.72*exp(-\t*\t/0.024)});
  \draw[blue!25, thin] (11.58,1.45) -- (12.42,1.45);

  \node[font=\footnotesize\itshape, text=black!55]
    at (6.0, 2.7) {$\mathcal{N}(\mu_k,\sigma_k)$ emission};

  % ===== Poisson clock (jump trigger box) =====
  \node[draw, rounded corners=5pt, fill=yellow!20, very thick,
        minimum width=2.8cm, minimum height=1.1cm,
        align=center, font=\small] (clock) at (6.0, 5.0)
        {Poisson clock\\[1pt]prob.~$\epsilon$};

  \draw[trigger] (sk.north) -- (clock.south)
    node[midway, right=3pt, font=\footnotesize] {jump triggered};

  % ===== Jump arcs from clock to tail state sets =====
  \draw[jumparc, red!70!black]
    (clock.west)
    to[out=180, in=100]
    node[above, sloped, font=\footnotesize, red!80!black,
         pos=0.45, yshift=2pt]
      {$K\!\sim\!\mathrm{Pois}(\lambda)$ forced steps}
    (s1.north);

  \draw[jumparc, blue!70!black]
    (clock.east)
    to[out=0, in=80]
    node[above, sloped, font=\footnotesize, blue!80!black,
         pos=0.45, yshift=2pt]
      {$K\!\sim\!\mathrm{Pois}(\lambda)$ forced steps}
    (sN.north);

  % ===== Tail set brace labels below states =====
  \draw[decorate,
        decoration={brace, amplitude=6pt, mirror},
        red!70!black, thick]
    (-0.6, -0.65) -- (2.6, -0.65)
    node[midway, below=7pt, red!80!black, font=\small]
      {$\mathcal{S}_{-}$\ (bottom tail)};

  \draw[decorate,
        decoration={brace, amplitude=6pt, mirror},
        blue!70!black, thick]
    (9.4, -0.65) -- (12.6, -0.65)
    node[midway, below=7pt, blue!80!black, font=\small]
      {$\mathcal{S}_{+}$\ (top tail)};

  % ===== State-axis arrow =====
  \draw[->, >=Stealth, gray!70, thin]
    (-1.0, -0.5) -- (13.2, -0.5)
    node[right, font=\small, text=black!70] {state index};

  \node[font=\footnotesize, text=gray!80, above right=2pt and 4pt of sk]
    {prob.~$1\!-\!\epsilon$};

\end{tikzpicture}%
}
\caption{Architecture of the Hybrid Hidden Markov Model with Jump-Diffusion (HMM-WJ).
    At each step the chain transitions according to the empirical transition matrix $\mathbf{T}$
    (probability $1-\epsilon$, thin bidirectional arrows) or enters a Poisson jump episode
    (probability $\epsilon$, dashed upward arrow to the Poisson clock).
    During a jump episode the active state is forced to the bottom tail set $\mathcal{S}_{-}$
    (red, lowest-valued states) or the top tail set $\mathcal{S}_{+}$ (blue, highest-valued states)
    for $K\!\sim\!\mathrm{Poisson}(\lambda)$ consecutive steps before normal Markovian evolution
    resumes.  Small bell curves above each state depict the state-conditional
    $\mathcal{N}(\mu_k,\sigma_k)$ emission distribution.  Tail sets are defined as the
    $N_{\rm tail}$ lowest- and highest-quantile states under the Laplace quantile partition.}
\label{fig:model_architecture}
\end{figure}


To characterize the excess growth rates within each regime, we parameterized the emission model with Normal distributions $\mathcal{N}(\mu_k, \sigma_k)$ whose bounds coincided with the 0.1 and 99.9 percentiles of the state-conditional distribution. Specifically, we set $\mu_k = (Q_{k-1} + Q_k)/2$ and $\sigma_k = (Q_k - Q_{k-1})/(2 z_{0.999})$, where $z_{0.999} \approx 3.09$ was the standard normal 99.9th percentile. During simulation, when the model was in state $k$, the continuous excess growth rate ($\hat{G}_t$) was sampled from the corresponding local distribution. This hybrid approach ensured that the model captured sudden large price moves while keeping smooth variation within each regime \cite{kou_jump-diffusion_2002}. The fitted Laplace CDF closely matched the empirical CDF for SPY, and the empirical transition matrix exhibited a dominant near-diagonal band reflecting short-range regime persistence (Figure~\ref{fig:model_internals}, Online Appendix~S1). Critically, under pure Markovian dynamics the model exited extreme (tail) states after only 1--2 steps on average, far shorter than the multi-week volatile episodes observed in real markets. This rapid reversion was the core limitation that motivated the jump-duration extension.

\subsection*{Growth rate generation.}
The basic sampling procedure for the HMM (without jumps) proceeded as follows. The stationary distribution $\bar{\pi}$, representing the long-run probability of occupying each regime, satisfies the balance equation $\bar{\pi} = \bar{\pi}\mathbf{T}$ and was estimated numerically by raising the transition matrix to a large power ($\mathbf{T}^{50}$) \cite{hamilton_regime_2018, grinstead_introduction_1997}. An initial hidden state was drawn from $X_0 \sim \bar{\pi}$, then for each subsequent step $n = 1, \dots, T$: (i)~the next hidden state was sampled from the transition row $X_n \sim \mathbf{T}_{X_{n-1},\,\cdot}$, and (ii)~an observation was sampled from the emission distribution $O_n \sim \mathbf{E}_{X_n}$. This produced a sequence of hidden states $\{X_1, \dots, X_T\}$ and corresponding observations $\{O_1, \dots, O_T\}$.

\subsection*{Hyperparameter optimization via grid search.}
To identify the optimal jump probability ($\epsilon$) and mean jump duration ($\lambda$) \cite{merton_jump_1976, kou_jump-diffusion_2002}, we implemented a multi-objective grid search that minimized the discrepancy between historical and simulated market signatures. The error function targeted two stylized facts \cite{cont_empirical_2001, ryden_stylized_1998}: volatility clustering, measured as the sum of squared errors between the observed and simulated autocorrelation function (ACF) of absolute excess growth rates \cite{schwert_why_1989} up to a lag of $L=252$ trading days (approximately one trading year); and leptokurtosis, captured by a penalty term on the global kurtosis \cite{mandelbrot_variation_1963}.

We defined the objective function $J(\epsilon, \lambda)$ as:
\begin{equation}
    J(\epsilon, \lambda) = \sum_{\tau=1}^{L} \left( \text{ACF}_{obs}(\tau) - \overline{\text{ACF}}_{sim}(\tau) \right)^2 + w_K \left( K_{obs} - \overline{K}_{sim} \right)^2
    \label{eq:grid_search}
\end{equation}
where $\text{ACF}_{obs}(\tau)$ and $K_{obs}$ represented the empirical absolute growth autocorrelation at lag $\tau$ and the global kurtosis, respectively. The simulated equivalents, $\overline{\text{ACF}}_{sim}(\tau)$ and $\overline{K}_{sim}$, were averaged across 200 independent synthetic paths of 2,766 trading days per grid point \cite{glasserman_monte_2003}. We set the kurtosis penalty weight to $w_K = 0.20$ to balance tail behavior against the temporal ACF term. The grid search explored $\epsilon$ values from $10^{-4}$ to $2.5\times10^{-2}$ and $\lambda$ values from $10$ to $160$. The full computational procedure is detailed in Algorithm \ref{alg:grid_search}.

\subsection*{Statistical validation metrics}
To assess whether synthetic paths were statistically consistent with empirical data, we applied two nonparametric goodness-of-fit tests to the marginal distributions of simulated excess growth rate sequences. For each of the 1,000 simulated paths, the empirical cumulative distribution function (CDF) of the simulated sequence was compared against the empirical CDF of the historical in-sample or out-of-sample observations.

The Kolmogorov-Smirnov (KS) statistic is defined as:
\begin{equation}
    D_{KS} = \sup_{x} \left\lvert F_n(x) - F_m(x) \right\rvert
    \label{eq:ks_stat}
\end{equation}
where $F_n$ and $F_m$ denote the empirical CDFs of the observed and simulated sequences of lengths $n$ and $m$, respectively \cite{kolmogorov_1933, smirnov_1948}. The null hypothesis of distributional equivalence is rejected at significance level $\alpha$ when $D_{KS}$ exceeds the critical value $c(\alpha)\sqrt{(n+m)/(nm)}$.

The Anderson-Darling (AD) statistic places greater weight on the distributional tails than the KS test and is defined as:
\begin{equation}
    A^2 = -n - \sum_{i=1}^{n} \frac{2i-1}{n}\left[\ln F(X_{(i)}) + \ln\left(1 - F(X_{(n+1-i)})\right)\right]
    \label{eq:ad_stat}
\end{equation}
where $X_{(1)} \le \cdots \le X_{(n)}$ are the order statistics of the pooled sample and $F$ is the empirical CDF of the reference distribution \cite{anderson_darling_1952}. The AD test is particularly suited for assessing tail-distributional fidelity, which is the primary concern for risk management applications.

We reported the \textit{pass rate} as the proportion of the 1,000 simulated paths for which the null hypothesis of distributional equivalence was not rejected at the $\alpha = 0.05$ significance level. A pass rate above 95 percent indicated that the synthetic paths were, in aggregate, statistically indistinguishable from the historical data under that test. Pass rates were reported separately for the in-sample and out-of-sample windows and for both model variants (HMM without jumps and the hybrid HMM with jumps).

\subsection*{Temporal fidelity metric}
While KS and AD tests assess marginal distributional quality, they are insensitive to temporal dependence. To measure how well a generator reproduced volatility clustering, we defined the autocorrelation mean absolute error (ACF-MAE) as:
\begin{equation}
    \text{ACF-MAE} = \frac{1}{L}\sum_{\tau=1}^{L} \left\lvert \widehat{\rho}_{|G|}^{\,\text{obs}}(\tau) - \widehat{\rho}_{|G|}^{\,\text{sim}}(\tau) \right\rvert
    \label{eq:acf_mae}
\end{equation}
where $\widehat{\rho}_{|G|}^{\,\text{obs}}(\tau)$ and $\widehat{\rho}_{|G|}^{\,\text{sim}}(\tau)$ are the sample autocorrelations of the absolute excess growth rate series $|G_t|$ at lag $\tau$ for the observed and simulated paths, respectively, and $L=252$ (one trading year). An i.i.d.\ generator produces $\widehat{\rho}_{|G|}^{\,\text{sim}}(\tau)\approx 0$ for all $\tau>0$, so its ACF-MAE equals the mean level of the observed autocorrelation function; lower values indicate better reproduction of empirical volatility persistence.

\subsection*{Single-Index Factor Model extension}
The framework described above applied to any individual asset ticker within the 424-asset dataset. To generalize synthetic path generation to a correlated multi-asset setting without fitting a separate HMM for each asset, we exploited the linear factor structure of the Single-Index Model (SIM) \cite{sharpe_simplified_1963}. The SIM expresses the excess growth rate of asset $i$ at time $t$ as:
\begin{equation}
    G_{i,t} = \alpha_i + \beta_i \cdot G_{\text{SPY},t} + \eta_{i,t}, \qquad \eta_{i,t} \stackrel{\text{i.i.d.}}{\sim} \mathcal{N}(0, \sigma_{\eta,i}^2)
    \label{eq:sim}
\end{equation}
where $G_{\text{SPY},t}$ denotes the SPY excess growth rate at time $t$, $\beta_i$ measures the systematic sensitivity of asset $i$ to the index, $\alpha_i$ is the idiosyncratic drift, and $\eta_{i,t}$ is a zero-mean idiosyncratic shock assumed independent across assets and time. The parameters $(\alpha_i, \beta_i, \sigma_{\eta,i}^2)$ were estimated via ordinary least squares on the in-sample time series for each asset $i$.

Synthetic multi-asset paths were generated in three steps. First, the hybrid HMM produced a single synthetic SPY excess growth rate path $\{\hat{G}_{\text{SPY},t}\}_{t=1}^T$ using the regime-switching and jump-duration mechanism described above. Second, for each asset $i$, an idiosyncratic shock sequence $\{\hat{\eta}_{i,t}\}$ was drawn independently from $\mathcal{N}(0, \hat{\sigma}_{\eta,i}^2)$. Third, the synthetic excess growth rate for asset $i$ was recovered as:
\begin{equation}
    \hat{G}_{i,t} = \hat{\alpha}_i + \hat{\beta}_i \cdot \hat{G}_{\text{SPY},t} + \hat{\eta}_{i,t}.
    \label{eq:sim_gen}
\end{equation}
This construction preserved the cross-sectional correlation structure induced by the common factor while maintaining the regime-switching and jump-diffusion dynamics of the index. The resulting 424-dimensional synthetic path could be used for stress testing, portfolio risk assessment, and scenario generation across the full asset universe in a single generative pass, avoiding the curse of dimensionality that would accompany a full multivariate HMM.

\subsection*{Simulation algorithm}
Algorithm~\ref{alg:sim} details the core jump-diffusion simulation procedure; the remaining pipeline stages (model construction, state decoding, and hyperparameter grid search) are provided in Online Appendix~S4.

\begin{algorithm}[tp]
\caption{Hybrid jump-diffusion simulation with persistent volatility}
\label{alg:sim}
\small
\begin{algorithmic}[1]
\REQUIRE Model parameters $\mathbf{T}, \bar{\pi}, \epsilon, \lambda, N_{tail}, p_{neg}$, Total Steps $M$
\ENSURE Simulated state sequence $S_{1:M}$
\STATE Define tail-states $\mathcal{S}_{bottom} = \{1, \dots, N_{tail}\}$ and $\mathcal{S}_{top} = \{N-N_{tail}+1, \dots, N\}$.
\STATE Sample initial state $S_1 \sim \text{Categorical}(\bar{\pi})$.
\STATE Set $counter \leftarrow 2$.
\WHILE{$counter \le M$}
    \STATE Sample $u \sim \text{Uniform}(0, 1)$.
    \IF{$u < \epsilon$}
        \STATE Sample jump-duration $K \sim \text{Poisson}(\lambda)$.
        \FOR{$j = 1$ to $K$ while $counter \le M$}
            \STATE Sample $w \sim \text{Uniform}(0, 1)$.
            \IF{$w < p_{neg}$}
                \STATE $S_{counter} \sim \text{Uniform}(\mathcal{S}_{bottom})$ \COMMENT{Bias toward negative-tail events}
            \ELSE
                \STATE $S_{counter} \sim \text{Uniform}(\mathcal{S}_{top})$
            \ENDIF
            \STATE $counter \leftarrow counter + 1$.
        \ENDFOR
    \ELSE
        \STATE Sample $S_{counter} \sim \text{Categorical}(\mathbf{T}_{S_{counter-1}, :})$.
        \STATE $counter \leftarrow counter + 1$.
    \ENDIF
\ENDWHILE
\RETURN $S_{1:M}$
\end{algorithmic}
\end{algorithm}


